%Version 3.1 December 2024
% See section 11 of the User Manual for version history
%
%%%%%%%%%%%%%%%%%%%%%%%%%%%%%%%%%%%%%%%%%%%%%%%%%%%%%%%%%%%%%%%%%%%%%%
%%                                                                 %%
%% Please do not use \input{...} to include other tex files.       %%
%% Submit your LaTeX manuscript as one .tex document.              %%
%%                                                                 %%
%% All additional figures and files should be attached             %%
%% separately and not embedded in the \TeX\ document itself.       %%
%%                                                                 %%
%%%%%%%%%%%%%%%%%%%%%%%%%%%%%%%%%%%%%%%%%%%%%%%%%%%%%%%%%%%%%%%%%%%%%

%%\documentclass[referee,sn-basic]{sn-jnl}% referee option is meant for double line spacing

%%=======================================================%%
%% to print line numbers in the margin use lineno option %%
%%=======================================================%%

%%\documentclass[lineno,pdflatex,sn-basic]{sn-jnl}% Basic Springer Nature Reference Style/Chemistry Reference Style

%%=========================================================================================%%
%% the documentclass is set to pdflatex as default. You can delete it if not appropriate.  %%
%%=========================================================================================%%

%%\documentclass[sn-basic]{sn-jnl}% Basic Springer Nature Reference Style/Chemistry Reference Style

%%Note: the following reference styles support Namedate and Numbered referencing. By default the style follows the most common style. To switch between the options you can add or remove “Numbered” in the optional parenthesis. 
%%The option is available for: sn-basic.bst, sn-chicago.bst%  
 
%%\documentclass[pdflatex,sn-nature]{sn-jnl}% Style for submissions to Nature Portfolio journals
%%\documentclass[pdflatex,sn-basic]{sn-jnl}% Basic Springer Nature Reference Style/Chemistry Reference Style
\documentclass[pdflatex,sn-mathphys-num]{sn-jnl}% Math and Physical Sciences Numbered Reference Style
%%\documentclass[pdflatex,sn-mathphys-ay]{sn-jnl}% Math and Physical Sciences Author Year Reference Style
%%\documentclass[pdflatex,sn-aps]{sn-jnl}% American Physical Society (APS) Reference Style
%%\documentclass[pdflatex,sn-vancouver-num]{sn-jnl}% Vancouver Numbered Reference Style
%%\documentclass[pdflatex,sn-vancouver-ay]{sn-jnl}% Vancouver Author Year Reference Style
%%\documentclass[pdflatex,sn-apa]{sn-jnl}% APA Reference Style
%%\documentclass[pdflatex,sn-chicago]{sn-jnl}% Chicago-based Humanities Reference Style

%%%% Standard Packages
%%<additional latex packages if required can be included here>

\usepackage{graphicx}%
\usepackage{multirow}%
\usepackage{amsmath,amssymb,amsfonts}%
\usepackage{amsthm}%
\usepackage{mathrsfs}%
\usepackage[title]{appendix}%
\usepackage{xcolor}%
\usepackage{textcomp}%
\usepackage{manyfoot}%
\usepackage{booktabs}%
\usepackage{algorithm}%
\usepackage{algorithmicx}%
\usepackage{algpseudocode}%
\usepackage{listings}%
%%%%
\usepackage{indentfirst}
\usepackage{epstopdf}
\usepackage{mathtools}
\usepackage{graphicx}
\usepackage{multirow}
\usepackage{amsmath,amssymb,amsfonts}
\usepackage{mathrsfs}
\usepackage{url}
\usepackage{xcolor}
\usepackage{textcomp}
\usepackage{manyfoot}
\usepackage{booktabs}
\usepackage{listings}
\usepackage{xspace}
\usepackage{rotating}
\usepackage{scalefnt}
\usepackage{upgreek}
\usepackage{enumitem}
\usepackage{array}
\usepackage{geometry}

\definecolor{keywordcolor}{rgb}{0.7, 0.1, 0.1}
\definecolor{tacticcolor}{rgb}{0.0, 0.1, 0.3}
\definecolor{commentcolor}{rgb}{0.4, 0.4, 0.4}
\definecolor{stringcolor}{rgb}{0.5, 0.3, 0.2}
\definecolor{symbolcolor}{rgb}{0.1, 0.2, 0.7}
\definecolor{sortcolor}{rgb}{0.1, 0.5, 0.1}
\definecolor{attributecolor}{rgb}{0.7, 0.1, 0.1}
\definecolor{errorcolor}{rgb}{1, 0, 0}

\lstdefinelanguage{lean}{
mathescape=false,
texcl=false,
morekeywords=[1]{
import, prelude, protected, private, noncomputable, definition, meta, renaming,
hiding, exposing, parameter, parameters, begin, conjecture, constant, constants,
hypothesis, lemma, corollary, variable, variables, premise, premises, theory,
print, theorem, proposition, example, open, as, export, override, axiom, axioms,
inductive, with, without, structure, record, universe, universes, alias, help,
precedence, reserve, declare_trace, add_key_equivalence, match, infix, infixl,
infixr, notation, postfix, prefix, instance, vm_eval, check, coercion, end,
this, suppose, using, using_well_founded, namespace, section, fields, attribute,
local, set_option, extends, include, omit, classes, class, instances, coercions,
attributes, raw, replacing, calc, have, show, suffices, by, in, at, let, forall,
Pi, fun, exists, if, dif, then, else, assume, take, obtain, from, aliases,
register_simp_ext, mutual, do, def, run_command, to, abbrev, where
},
morekeywords=[2]{Sort, Type, Prop, Type*, Type₀, Type₁, Type₂, Type₃},
morekeywords=[3]{sorry, admit},
literate=
{α}{{\ensuremath{\mathrm{\upalpha}}}}1
{β}{{\ensuremath{\mathrm{\upbeta}}}}1
{γ}{{\ensuremath{\mathrm{\upgamma}}}}1
{δ}{{\ensuremath{\mathrm{\updelta}}}}1
{ε}{{\ensuremath{\mathrm{\varepsilon}}}}1
{ζ}{{\ensuremath{\mathrm{\zeta}}}}1
{η}{{\ensuremath{\mathrm{\eta}}}}1
{θ}{{\ensuremath{\mathrm{\theta}}}}1
{ι}{{\ensuremath{\mathrm{\iota}}}}1
{κ}{{\ensuremath{\mathrm{\kappa}}}}1
{λ}{{\color{symbolcolor}\ensuremath{\mathrm{\uplambda}}}}1
{μ}{{\ensuremath{\mathrm{\mu}}}}1
{ν}{{\ensuremath{\mathrm{\nu}}}}1
{ξ}{{\ensuremath{\mathrm{\xi}}}}1
{π}{{\ensuremath{\mathrm{\mathnormal{\pi}}}}}1
{ρ}{{\ensuremath{\mathrm{\rho}}}}1
{σ}{{\ensuremath{\mathrm{\sigma}}}}1
{τ}{{\ensuremath{\mathrm{\tau}}}}1
{χ}{{\ensuremath{\mathrm{\chi}}}}1
{ψ}{{\ensuremath{\mathrm{\psi}}}}1
{ω}{{\ensuremath{\mathrm{\omega}}}}1
{υ}{{\ensuremath{\mathrm{\nu}}}}1
{Γ}{{\ensuremath{\mathrm{\Gamma}}}}1
{Δ}{{\ensuremath{\mathrm{\Delta}}}}1
{Θ}{{\ensuremath{\mathrm{\Theta}}}}1
{Λ}{{\ensuremath{\mathrm{\Lambda}}}}1
{Σ'}{{\color{symbolcolor}\ensuremath{\Sigma'}}}2
{Σ}{{\color{symbolcolor}\ensuremath{\Sigma}}}1
{Ξ}{{\ensuremath{\mathrm{\Xi}}}}1
{Ψ}{{\ensuremath{\mathrm{\Psi}}}}1
{Ω}{{\ensuremath{\mathrm{\Omega}}}}1
{Φ}{{\ensuremath{\mathrm{\Phi}}}}1
{Π}{{\color{symbolcolor}\ensuremath{\mathrm{\Uppi}}}}1
{ℵ}{{\ensuremath{\aleph}}}1
{ℕ}{{\ensuremath{\mathbb{N}}}}1
{ℤ}{{\ensuremath{\mathbb{Z}}}}1
{ℝ}{{\ensuremath{\mathbb{R}}}}1
{ℚ}{{\ensuremath{\mathbb{Q}}}}1
{ℂ}{{\ensuremath{\mathbb{C}}}}1
{ℓ}{{\ensuremath{\ell}}}1
{⊤}{{\ensuremath{\top}}}1
{⊥}{{\color{symbolcolor}\ensuremath{\perp}}}1
{∞}{{\color{symbolcolor}\ensuremath{\infty}}}1
{√}{{\color{symbolcolor}\ensuremath{\sqrt}}}1
{≤}{{\color{symbolcolor}\ensuremath{\leq}}}1
{≥}{{\color{symbolcolor}\ensuremath{\geq}}}1
{≠}{{\color{symbolcolor}\ensuremath{\neq}}}1
{≈}{{\color{symbolcolor}\ensuremath{\approx}}}1
{≡}{{\color{symbolcolor}\ensuremath{\equiv}}}1
{≃o}{{\color{symbolcolor}\ensuremath{\simeq_o}}}2
{≃}{{\color{symbolcolor}\ensuremath{\simeq}}}1
{∂}{{\color{symbolcolor}\ensuremath{\partial}}}1
{∆}{{\color{symbolcolor}\ensuremath{\triangle}}}1
{∫}{{\color{symbolcolor}\ensuremath{\int}}}1
{∑}{{\color{symbolcolor}\ensuremath{\Sigma}}}1
{∓}{{\color{symbolcolor}\ensuremath{\mp}}}1
{±}{{\color{symbolcolor}\ensuremath{\pm}}}1
{×}{{\color{symbolcolor}\ensuremath{\times}}}1
{⊕ₗ}{{\color{symbolcolor}\ensuremath{\oplus_{\ell}}}}2
{⊕}{{\color{symbolcolor}\ensuremath{\oplus}}}1
{⊗}{{\color{symbolcolor}\ensuremath{\otimes}}}1
{⊞}{{\color{symbolcolor}\ensuremath{\boxplus}}}1
{⨆}{{\ensuremath{\bigsqcup}}}1
{⨅}{{\color{symbolcolor}\ensuremath{\bigsqcap}}}1
{∇}{{\color{symbolcolor}\ensuremath{\nabla}}}1
{⬝}{{\color{symbolcolor}\ensuremath{\cdot}}}1
{•}{{\color{symbolcolor}\ensuremath{\cdot}}}1
{∘}{{\color{symbolcolor}\ensuremath{\circ}}}1
{∧}{{\color{symbolcolor}\ensuremath{\wedge}}}1
{∨}{{\color{symbolcolor}\ensuremath{\vee}}}1
{¬}{{\color{symbolcolor}\ensuremath{\neg}}}1
{⊢}{{\color{symbolcolor}\ensuremath{\vdash}}}1
{∀}{{\color{symbolcolor}\ensuremath{\forall}}}1
{∃}{{\color{symbolcolor}\ensuremath{\exists}}}1
{↦}{{\color{symbolcolor}\ensuremath{\mapsto}}}1
{→}{{\color{symbolcolor}\ensuremath{\rightarrow}}}1
{←}{{\color{symbolcolor}\ensuremath{\leftarrow}}}1
{↔}{{\color{symbolcolor}\ensuremath{\leftrightarrow}}}1
{⇒}{{\color{symbolcolor}\ensuremath{\Rightarrow}}}1
{⟹}{{\color{symbolcolor}\ensuremath{\Longrightarrow}}}1
{⇐}{{\color{symbolcolor}\ensuremath{\Leftarrow}}}1
{⟸}{{\color{symbolcolor}\ensuremath{\Longleftarrow}}}1
{⟶}{{\color{symbolcolor}\ensuremath{\longrightarrow}}}1
{⥤}{{\color{symbolcolor}\ensuremath{\Rightarrow}}}1
{∩}{{\color{symbolcolor}\ensuremath{\cap}}}1
{∪}{{\color{symbolcolor}\ensuremath{\cup}}}1
{⊂}{{\color{symbolcolor}\ensuremath{\subseteq}}}1
{⊆}{{\color{symbolcolor}\ensuremath{\subseteq}}}1
{⊄}{{\color{symbolcolor}\ensuremath{\nsubseteq}}}1
{⊈}{{\color{symbolcolor}\ensuremath{\nsubseteq}}}1
{⊃}{{\color{symbolcolor}\ensuremath{\supseteq}}}1
{⊇}{{\color{symbolcolor}\ensuremath{\supseteq}}}1
{⊅}{{\color{symbolcolor}\ensuremath{\nsupseteq}}}1
{⊉}{{\color{symbolcolor}\ensuremath{\nsupseteq}}}1
{∈}{{\color{symbolcolor}\ensuremath{\in}}}1
{∉}{{\color{symbolcolor}\ensuremath{\notin}}}1
{∋}{{\color{symbolcolor}\ensuremath{\ni}}}1
{∌}{{\color{symbolcolor}\ensuremath{\notni}}}1
{∅}{{\color{symbolcolor}\ensuremath{\emptyset}}}1
{∖}{{\color{symbolcolor}\ensuremath{\setminus}}}1
{†}{{\color{symbolcolor}\ensuremath{\dag}}}1
{⋆}{{\ensuremath{\star}}}1
{∥}{{\ensuremath{\|}}}1
{φ}{{\ensuremath{\phi}}}1
{⌞}{{\ensuremath{\llcorner}}}1
{⌟}{{\ensuremath{\lrcorner}}}1
{⦃}{{\ensuremath{\{\!|}}}1
{⦄}{{\ensuremath{|\!\}}}}1
{⟨}{{\ensuremath{\langle}}}1
{⟩}{{\ensuremath{\rangle}}}1
{⟪}{{\ensuremath{\llangle}}}1
{⟫}{{\ensuremath{\rrangle}}}1
{⌊}{{\ensuremath{\lfloor}}}1
{⌋}{{\ensuremath{\rfloor}}}1
{⌈}{{\ensuremath{\lceil}}}1
{⌉}{{\ensuremath{\rceil}}}1
{‹}{{\guilsinglleft}}1
{›}{{\guilsinglright}}1
{｜}{{\ensuremath{\mid}}}1
{₀₀}{{\ensuremath{_{00}}}}2
{₁₁}{{\ensuremath{_{11}}}}2
{₁₂}{{\ensuremath{_{12}}}}2
{₂₁}{{\ensuremath{_{21}}}}2
{₂₂}{{\ensuremath{_{22}}}}2
{₀}{{\ensuremath{_0}}}1
{₁}{{\ensuremath{_1}}}1
{₂}{{\ensuremath{_2}}}1
{₃}{{\ensuremath{_3}}}1
{₄}{{\ensuremath{_4}}}1
{₅}{{\ensuremath{_5}}}1
{₆}{{\ensuremath{_6}}}1
{₇}{{\ensuremath{_7}}}1
{₈}{{\ensuremath{_8}}}1
{₉}{{\ensuremath{_9}}}1
{₊}{{\ensuremath{_+}}}1
{ₐ}{{\ensuremath{_a}}}1
{ᵢ}{{\ensuremath{_i}}}1
{ⱼ}{{\ensuremath{_j}}}1
{ₗ}{{\ensuremath{_l}}}1
{ₙ}{{\ensuremath{_n}}}1
{ₘ}{{\ensuremath{_m}}}1
{ₛ}{{\ensuremath{_s}}}1
{ᵥ}{{\ensuremath{_v}}}1
{ᵀ}{{\ensuremath{^\top}}}1
{ᵒᵖ}{{\color{symbolcolor}\textsuperscript{op}}}2
{ᗮ}{{\ensuremath{^\perp}}}1
{¹}{{\ensuremath{^1}}}1
{ᶠ}{{\ensuremath{^f}}}1
{ˡ}{{\ensuremath{^l}}}1
{ⁱ}{{\ensuremath{^i}}}1
{⁻}{{\ensuremath{^{-}}}}1
{↑}{{\color{symbolcolor}\ensuremath{\uparrow}}}1
{↓}{{\color{symbolcolor}\ensuremath{\downarrow}}}1
{//}{{\color{symbolcolor}//}}2
{<|>}{{\color{symbolcolor}<|>}}3
{...}{{\ensuremath{\ldots}}}3
{\$}{{\color{symbolcolor}\$}}1
{:}{{\color{symbolcolor}:}}1
{|}{{\color{symbolcolor}|}}1
{=}{{\color{symbolcolor}=}}1
{<}{{\color{symbolcolor}<}}1
{>}{{\color{symbolcolor}>}}1
{+}{{\color{symbolcolor}+}}1
{*}{{\color{symbolcolor}*}}1
{`}{{\ensuremath{{}^\backprime}}}1
{'}{{\ensuremath{{}^\prime}}}1,
morecomment=[s]{/-}{-/},
morecomment=[l]{--},
showstringspaces=false,
keepspaces=true,
morestring=[b]{"},
tabsize=3,
extendedchars=false,
sensitive=true,
breaklines=true,
breakatwhitespace=true,
lineskip={-1.5pt},
basicstyle={\ttfamily\scalefont{.95}},
captionpos=b,
columns=[l]fullflexible,
identifierstyle={\ttfamily\color{black}},
keywordstyle=[1]{\ttfamily\color{keywordcolor}},
keywordstyle=[2]{\ttfamily\color{sortcolor}},
keywordstyle=[3]{\ttfamily\color{errorcolor}},
stringstyle={\ttfamily\color{stringcolor}},
commentstyle={\ttfamily\itshape\color{commentcolor}},
}

\newcommand{\lean}[1]{\lstinline[language=lean, mathescape=true]{#1}}
\lstset{language=lean}
\newcommand{\re}[1]{\textcolor{red}{#1}}
\newcommand{\mathlib}{\textsf{mathlib}\xspace}



%%%%%=============================================================================%%%%
%%%%  Remarks: This template is provided to aid authors with the preparation
%%%%  of original research articles intended for submission to journals published 
%%%%  by Springer Nature. The guidance has been prepared in partnership with 
%%%%  production teams to conform to Springer Nature technical requirements. 
%%%%  Editorial and presentation requirements differ among journal portfolios and 
%%%%  research disciplines. You may find sections in this template are irrelevant 
%%%%  to your work and are empowered to omit any such section if allowed by the 
%%%%  journal you intend to submit to. The submission guidelines and policies 
%%%%  of the journal take precedence. A detailed User Manual is available in the 
%%%%  template package for technical guidance.
%%%%%=============================================================================%%%%

%% as per the requirement new theorem styles can be included as shown below
\theoremstyle{thmstyleone}%
\newtheorem{theorem}{Theorem}%  meant for continuous numbers
%%\newtheorem{theorem}{Theorem}[section]% meant for sectionwise numbers
%% optional argument [theorem] produces theorem numbering sequence instead of independent numbers for Proposition
\newtheorem{proposition}[theorem]{Proposition}% 
\newtheorem{lemma}[theorem]{Lemma}
%%\newtheorem{proposition}{Proposition}% to get separate numbers for theorem and proposition etc.

\theoremstyle{thmstyletwo}%
\newtheorem{example}{Example}%
\newtheorem{remark}{Remark}%

\theoremstyle{thmstylethree}%
\newtheorem{definition}{Definition}%

\raggedbottom
%%\unnumbered% uncomment this for unnumbered level heads

\begin{document}

\title[Unified Framework for Matrix Decompositions]{A Unified Formal Framework for Existence Proofs of Matrix Decompositions}

%%=============================================================%%
%% GivenName	-> \fnm{Joergen W.}
%% Particle	-> \spfx{van der} -> surname prefix
%% FamilyName	-> \sur{Ploeg}
%% Suffix	-> \sfx{IV}
%% \author*[1,2]{\fnm{Joergen W.} \spfx{van der} \sur{Ploeg} 
%%  \sfx{IV}}\email{iauthor@gmail.com}
%%=============================================================%%

\author*[1,2]{\fnm{First} \sur{Author}}\email{iauthor@gmail.com}

\author[2,3]{\fnm{Second} \sur{Author}}\email{iiauthor@gmail.com}
\equalcont{These authors contributed equally to this work.}

\author[1,2]{\fnm{Third} \sur{Author}}\email{iiiauthor@gmail.com}
\equalcont{These authors contributed equally to this work.}

\affil*[1]{\orgdiv{Department}, \orgname{Organization}, \orgaddress{\street{Street}, \city{City}, \postcode{100190}, \state{State}, \country{Country}}}

\affil[2]{\orgdiv{Department}, \orgname{Organization}, \orgaddress{\street{Street}, \city{City}, \postcode{10587}, \state{State}, \country{Country}}}

\affil[3]{\orgdiv{Department}, \orgname{Organization}, \orgaddress{\street{Street}, \city{City}, \postcode{610101}, \state{State}, \country{Country}}}

%%==================================%%
%% Sample for unstructured abstract %%
%%==================================%%

\abstract{
Matrix decompositions are fundamental in linear algebra and scientific computing, and their existence proofs often follow a common pattern: a matrix is transformed into a suitable form, reduced to a smaller subproblem, solved recursively, and then reconstructed through block arguments. This paper proposes a unified formal framework for such proofs in Lean 4. The framework separates decomposition goals, local transformations, reduction methods, branching strategies, and the global induction mechanism into modular components. To support recursive arguments across changing matrix dimensions, we introduce a matrix universe together with a well-founded induction engine. We also develop reusable infrastructure for block lifting, reindexing, and the transfer of structural properties. The framework is validated through two case studies. The first is a detailed formalization of the existence of the lower-upper decomposition with partial pivoting (PLU), which serves as the main example. The second is a structurally different decomposition instance, used to show that the framework is reusable beyond a single specialized setting. These results indicate that existence proofs for matrix decompositions can be organized around reusable proof patterns, and they provide a basis for further developments in formal mathematics for structured linear algebra.
}

%%================================%%
%% Sample for structured abstract %%
%%================================%%

% \abstract{\textbf{Purpose:} The abstract serves both as a general introduction to the topic and as a brief, non-technical summary of the main results and their implications. The abstract must not include subheadings (unless expressly permitted in the journal's Instructions to Authors), equations or citations. As a guide the abstract should not exceed 200 words. Most journals do not set a hard limit however authors are advised to check the author instructions for the journal they are submitting to.
% 
% \textbf{Methods:} The abstract serves both as a general introduction to the topic and as a brief, non-technical summary of the main results and their implications. The abstract must not include subheadings (unless expressly permitted in the journal's Instructions to Authors), equations or citations. As a guide the abstract should not exceed 200 words. Most journals do not set a hard limit however authors are advised to check the author instructions for the journal they are submitting to.
% 
% \textbf{Results:} The abstract serves both as a general introduction to the topic and as a brief, non-technical summary of the main results and their implications. The abstract must not include subheadings (unless expressly permitted in the journal's Instructions to Authors), equations or citations. As a guide the abstract should not exceed 200 words. Most journals do not set a hard limit however authors are advised to check the author instructions for the journal they are submitting to.
% 
% \textbf{Conclusion:} The abstract serves both as a general introduction to the topic and as a brief, non-technical summary of the main results and their implications. The abstract must not include subheadings (unless expressly permitted in the journal's Instructions to Authors), equations or citations. As a guide the abstract should not exceed 200 words. Most journals do not set a hard limit however authors are advised to check the author instructions for the journal they are submitting to.}

\keywords{Matrix Decomposition, Existence Proof, Linear Algebra, Well-founded Induction, Block Lifting; Formal Mathematics}

%%\pacs[JEL Classification]{D8, H51}

%%\pacs[MSC Classification]{35A01, 65L10, 65L12, 65L20, 65L70}

\maketitle
% 小标题不要超过一行
\section{Introduction}
\label{sec:introduction}

% \subsection{(Background)}
Matrix decompositions are fundamental in linear algebra and scientific computing. They underlie a wide range of tasks, including the solution of linear systems, least-squares problems, eigenvalue computations, and low-rank approximation, and they provide the structural basis for many classical numerical algorithms \cite{stewart1998matrix,trefethen2022numerical}. Meanwhile, formalization platforms such as Lean and Mathlib have developed substantial foundations for algebra, analysis, topology, and linear algebra \cite{mathlib2020,avigad2020mathematics}. This makes matrix decompositions a natural and important class of targets for formalized mathematics.

% \subsection{(Problem Statement and Related Work)}
% \

Although matrix decompositions take different concrete forms, their existence proofs often follow a similar general pattern: one first applies suitable local transformations to the original matrix, reduces the problem to a smaller subproblem, establishes the desired result recursively for that subproblem, and finally lifts the result back to the original matrix by means of a block construction. Similar proof structures appear in matrix reduction procedures, constructions of normal forms, and related results in linear algebra \cite{aransay2016formalisation,divason2022formalization}. This suggests that different matrix decompositions are connected not only by their subject matter, but also by structural commonalities in their proofs.

Substantial foundations for formalized work in linear algebra and matrix theory have already been established. Formalization platforms represented by Lean and its mathematical library mathlib provide a unified foundational environment for algebra, analysis, topology, and linear algebra, while the improvements introduced in Lean 4 further strengthen the feasibility of modular formal developments through its language mechanisms, metaprogramming facilities, and system extensibility \cite{mathlib2020,moura2021lean}. In work closer to the present topic, Aransay and Divasón formalized the computation of matrix echelon forms in Isabelle/HOL and proved the correctness of the corresponding algorithm \cite{aransay2016formalisation}. Divasón and Thiemann later developed a systematic formalization of the Smith normal form, proving the correctness of the relevant algorithm in a more general algebraic setting \cite{divason2022formalization}. These results show that classical statements involving matrix transformations, reduction procedures, and canonical forms can already be expressed and verified rigorously in existing proof assistants \cite{aransay2016formalisation,divason2022formalization}.

At the same time, most existing developments are organized around a single algorithm, a single normal form, or an individual theorem, and their proof structures are therefore closely tied to the specific target under consideration \cite{aransay2016formalisation,divason2022formalization}. By contrast, several recurring components of existence proofs for matrix decompositions—such as local transformations, reductions to subproblems, cross-dimensional recursion, and the recovery of a global result from a solved subproblem—still lack a unified formal organization. The present paper is concerned with this level of structure: how to elevate these shared proof components into a reusable formal framework and use it to support systematic constructions for different matrix decomposition instances.

% \subsection{(Main Contributions)}
% \

This paper develops a unified formal framework for existence proofs of matrix decompositions and validates it through two representative case studies. The main contributions of the paper can be summarized as follows:
\begin{enumerate}
% it改成we
    \item We proposes a unified formal framework for existence proofs of matrix decompositions, in which decomposition goals, local transformations, reduction methods, branching strategies, and the global induction mechanism are organized in a layered and coherent manner, thereby providing a common proof interface for different decomposition instances.

    \item We introduces a formal method for cross-dimensional matrix recursion based on a matrix universe and a well-founded induction mechanism, which systematically connects lower-dimensional subproblems with global existence proofs and resolves the difficulty that recursive constructions involve matrices of varying dimensions.

    \item We develops a collection of reusable formal components, including infrastructure for block lifting, reindexing, and the transfer of structural properties under embeddings and transformations, thereby systematizing several technical steps that would otherwise be repeated in individual formalizations.

    \item We completes two representative case studies: the first uses the PLU decomposition as the main example and exhibits in full how the different layers of the framework and its supporting components interact; the second considers a decomposition of substantially different structural character, showing that the framework is not tailored to a single decomposition and has potential for broader applicability.
\end{enumerate}

The emphasis of this work is not on formalizing one isolated decomposition theorem in isolation, but on developing a reusable organizational method for a class of results that share a common proof structure. The proposed framework and its case studies show that several core mechanisms in existence proofs of matrix decompositions can be expressed in a uniform way, thereby providing a formal basis for structured proofs in linear algebra.

% \subsection{(Organization of the Paper)}
% \

The remainder of the paper is organized as follows. Section 2 introduces the basic formal background and summarizes the common mathematical pattern underlying existence proofs of matrix decompositions, together with the main difficulties that arise in their formalization. Section 3 presents the unified formal framework, explaining how decomposition goals, local transformations, reduction methods, and branching strategies are organized. Section 4 develops the core technical machinery and reusable components supporting the framework. Section 5 takes the PLU decomposition as the main case study and presents its concrete instantiation within the framework together with the main formalization results. Section 6 discusses a second decomposition instance, showing how the framework can be reused and adapted under a different target structure. Section 7 examines the extensibility, limitations, and methodological significance of the framework. Finally, Section 8 concludes the paper and outlines directions for future work.

\section{Preliminaries and Motivation}
\label{sec:preliminaries}
% 强调一下我们要解决的问题
\subsection{Matrices in Lean 4 and Mathlib}
\

The formal developments in this paper are based on Lean 4 and its mathematical library mathlib~\cite{mathlib2020,moura2021lean}. In mathlib, a matrix is represented as a function on index types:
\[
\mathrm{Matrix}\; m\; n\; \alpha \;\equiv\; m \to n \to \alpha .
\]
Thus, for finite-dimensional matrices, an $n \times n$ matrix is typically written as \texttt{Matrix (Fin n) (Fin n) $\alpha$}. This representation places dimension information directly in the type, so that many shape constraints are enforced by typing rather than added as separate assumptions~\cite{mathlib2020,avigad2020mathematics}.

This design is well suited to formal linear algebra, but it also creates a characteristic difficulty. Passing from an $n \times n$ matrix to an $(n-1) \times (n-1)$ submatrix is not merely a change of notation: it means moving to a different matrix type. Similarly, block constructions, row and column permutations, and transfers between finite index types must be handled through explicit reindexing maps and equivalences. As a result, recursive arguments across changing dimensions require additional formal infrastructure.

Mathlib already provides a substantial collection of matrix operations and basic constructions, including multiplication, transpose, submatrices, and reindexing under index equivalences~\cite{mathlib2020,avigad2020mathematics}. These foundations make it possible to formulate matrix decomposition problems in a uniform language, while leaving the main challenge at the level of proof organization. This is precisely the setting in which the framework developed in the following sections is situated.

\subsection{A Common Mathematical Pattern in Existence Proofs}
% 介绍几个常见的矩阵分解
The existence proof of the Cholesky decomposition is one of the clearest places to begin. Let \(A \in \mathbb{R}^{n \times n}\) be symmetric positive definite, and write it in block form as
\[
A=
\begin{pmatrix}
a & b^{T}\\
b & B
\end{pmatrix}.
\]
Since \(A\) is positive definite, one has \(a>0\). One may therefore define
\[
r_{11}=\sqrt a,
\qquad
w=\frac{1}{r_{11}}\,b,
\]
and consider the matrix
\[
S=B-ww^{T}=B-\frac{1}{a}bb^{T}.
\]
A standard result shows that \(S\) is again symmetric positive definite. The original problem is thus reduced to the same kind of problem for the smaller matrix \(S\). If recursively one already has
\[
S=R_1^{T}R_1,
\]
then defining
\[
R=
\begin{pmatrix}
r_{11} & w^{T}\\
0 & R_1
\end{pmatrix}
\]
gives
\[
A=R^{T}R.
\]
This proof does not construct the whole factor \(R\) at once. Instead, it first extracts a part determined by the leading entry and column, compresses the remaining information into the Schur complement \(S\), solves the same problem recursively for \(S\), and finally restores the result to the original dimension by a block matrix construction.

The existence proof of an LU or PLU decomposition has an obvious resemblance to this, although the first step now takes the form of an explicit transformation. Let \(A \in K^{n \times n}\). If the first column is zero, then \(A\) can be written as
\[
A=
\begin{pmatrix}
0 & *\\
0 & A'
\end{pmatrix},
\]
so the problem immediately reduces to \(A' \in K^{(n-1)\times(n-1)}\). If the first column contains a nonzero entry, one chooses a pivot and moves it to the upper-left corner by a permutation. After this step one may assume that
\[
A=
\begin{pmatrix}
a_{11} & r\\
c & B
\end{pmatrix},
\qquad a_{11}\neq 0.
\]
Eliminating below the first entry yields
\[
L_1^{-1}A=
\begin{pmatrix}
a_{11} & r\\
0 & B-ca_{11}^{-1}r
\end{pmatrix}.
\]
The lower-right block
\[
S=B-ca_{11}^{-1}r
\]
is then the new subproblem. Once a decomposition of \(S\) is obtained recursively, the smaller factors can be embedded into the lower-right block and combined with the earlier permutation and elimination steps to recover a decomposition of the original matrix. Compared with Cholesky, this proof features pivot selection, row exchange, and elimination as explicit local operations, but its overall progression is similar: a local step is performed first, the problem is shifted to a smaller trailing block, and the final result is assembled from the local transformation and the recursive solution.

QR decomposition provides a rather different example. Let \(A\in\mathbb{R}^{m\times n}\), and let \(a_1\) denote its first column. A Householder reflection produces an orthogonal matrix \(H_1\) such that
\[
H_1a_1=
\begin{pmatrix}
\pm\|a_1\|\\
0\\
\vdots\\
0
\end{pmatrix}.
\]
Hence
\[
H_1A=
\begin{pmatrix}
* & *\\
0 & A_1
\end{pmatrix},
\]
where \(A_1\) is a smaller trailing submatrix. If recursively one has
\[
A_1=Q_1R_1,
\]
then one may define
\[
\widetilde Q=
\begin{pmatrix}
1 & 0\\
0 & Q_1
\end{pmatrix},
\qquad
\widetilde R=
\begin{pmatrix}
* & *\\
0 & R_1
\end{pmatrix},
\]
so that
\[
H_1A=\widetilde Q\,\widetilde R,
\qquad
A=H_1^{T}\widetilde Q\,\widetilde R.
\]
Since \(H_1^{T}\widetilde Q\) is still orthogonal, this yields the desired QR decomposition. The algebraic objects here are quite different from those in LU or Cholesky: there is no pivoting, and the central tool is an orthogonal reflection. Yet the proof still proceeds by first performing a local operation that exposes a suitable block structure, then passing to a smaller subproblem of the same type, and finally lifting the result back to the original dimension.

A different kind of example is provided by rank-normal-form arguments. Let \(A \in K^{m\times n}\) with \(\operatorname{rank}(A)=r\). A classical theorem states that there exist invertible matrices \(P\) and \(Q\) such that
\[
PAQ=
\begin{pmatrix}
I_r & 0\\
0 & 0
\end{pmatrix}.
\]
Its proof typically begins with a nonzero entry of \(A\). By row and column exchanges this entry is moved to the upper-left corner, and by elementary operations on both sides the first row and column are brought into a form such as
\[
\begin{pmatrix}
1 & 0\\
0 & *
\end{pmatrix}.
\]
This leaves a smaller lower-right block \(A_1\). The same procedure is then applied recursively to \(A_1\), and the row and column operations from all stages are accumulated to produce the full normal-form representation. Here the relevant tools are no longer triangular factors or orthogonal matrices, but two-sided elementary transformations. Even so, the proof is again built from the same kind of progression: a local normalization is performed first, a smaller remainder is identified, that remainder is handled recursively, and the local steps are finally combined into a global result.

Placed side by side, these proofs suggest a pattern that is not visible at the level of formulas alone. Each proof begins with a local operation on the current matrix: in Cholesky one extracts the part determined by the leading block, in LU or PLU one performs pivoting and elimination, in QR one applies an orthogonal reflection, and in rank normal form one uses elementary operations on both sides. After this local step, the matrix exhibits a block structure, and the main task is transferred to a smaller object. That smaller object may be a Schur complement, a trailing submatrix, or a reduced block left after normalization. It is typically of the same general type as the original problem, so the same decomposition theorem can be applied recursively. Once that recursive step is completed, the result must still be lifted back to the original matrix and combined with the preceding local operations before the global decomposition is obtained.

It is therefore natural to distinguish four recurrent kinds of mathematical operations in these proofs: a \emph{local transformation}, a \emph{reduction to a smaller subproblem}, a \emph{recursive construction}, and a \emph{reconstruction of the global result}. This terminology is not meant to impose an artificial template on all matrix decompositions. Rather, it is a way of naming the recurring proof steps that become visible when one compares classical existence arguments of rather different kinds.

\begin{figure}[t]
  \centering
  % figure to be inserted
  \includegraphics[width=\linewidth]{fig/proof-patern-figure.pdf}
  \caption{A schematic proof pattern for matrix decomposition existence theorems.}
  \label{fig:proof-pattern}
\end{figure}

% Figure~\ref{fig:proof-pattern} summarizes this observation in an abstract form. The figure does not correspond to any single decomposition theorem. Instead, it distills a recurrent proof organization from the examples above: a local operation first exposes a block structure containing a smaller subproblem \(A_1\); the same kind of decomposition is then constructed recursively for \(A_1\); and the resulting factors are lifted back to the original dimension to recover a global decomposition of \(A\). What varies from one decomposition to another is the concrete nature of the local operation, the precise shape of the reduced subproblem, and the structural properties that must be preserved during reconstruction. The underlying proof progression is nevertheless strikingly stable across these representative cases. The formal framework developed in the following sections is designed precisely to capture this level of common structure.

\begin{theorem}[An abstract induction principle on a controlled subtype]
\label{thm:abstract-subtype-induction}
Let $X$ be a set, let $\iota : S \to X$ be an embedding of a controlled subtype $S$ into $X$, and let
\[
\mu : X \to \alpha
\]
be a measure into a set $\alpha$ equipped with a well-founded relation $\prec$. Let $P$ be a property on $X$, and let $P_S$ be a property on $S$, such that
\[
P_S(s)\quad \Longleftrightarrow \quad P(\iota(s))
\qquad
\text{for all } s \in S.
\]
Assume further that:
\begin{enumerate}
  \item there is a distinguished set of base objects $B \subseteq X$ such that
  \[
  x \notin \iota(S)\ \text{or}\ x \in B
  \quad \Longrightarrow \quad
  P(x)
  \qquad
  \text{for all } x \in X;
  \]
  \item there is a relation $\leadsto$ on $S$ such that
  \[
  s' \leadsto s \ \text{and}\ P_S(s')
  \quad \Longrightarrow \quad
  P_S(s)
  \qquad
  \text{for all } s',s \in S;
  \]
  \item there is a class of sliceable objects in $S$, together with a slicing operation
  \[
  s \longmapsto \operatorname{slice}(s) \in X,
  \]
  such that whenever $s \in S$ is sliceable and
  \[
  P(\operatorname{slice}(s))
  \]
  holds, then
  \[
  P_S(s)
  \]
  holds;
  \item for every $s \in S$ with $\iota(s) \notin B$, there exists $t \in S$ such that
  \begin{enumerate}
    \item $t \leadsto s$,
    \item $t$ is sliceable, and
    \item
    \[
    \mu(\operatorname{slice}(t)) \prec \mu(\iota(s)).
    \]
  \end{enumerate}
\end{enumerate}
Then
\[
\forall x \in X,\quad P(x).
\]
\end{theorem}

\begin{proof}
We argue by well-founded induction on the value of $\mu(x)$. Fix $x \in X$, and assume as induction hypothesis that
\[
P(y)
\qquad
\text{holds for all } y \in X \text{ with } \mu(y) \prec \mu(x).
\]
We show that $P(x)$ holds.

If $x \notin \iota(S)$ or $x \in B$, the claim follows immediately from assumption~(1). It remains to consider the case in which $x$ lies in the image of $\iota$ and $x \notin B$. Since $\iota$ is an embedding, there exists $s \in S$ such that
\[
x = \iota(s).
\]
Because $\iota(s) \notin B$, assumption~(4) yields an element $t \in S$ such that
\[
t \leadsto s,
\qquad
t \text{ is sliceable},
\qquad
\mu(\operatorname{slice}(t)) \prec \mu(\iota(s)).
\]
Since $x=\iota(s)$, the last inequality becomes
\[
\mu(\operatorname{slice}(t)) \prec \mu(x).
\]
The induction hypothesis therefore applies to $\operatorname{slice}(t)$ and gives
\[
P(\operatorname{slice}(t)).
\]
Because $t$ is sliceable, assumption~(3) implies
\[
P_S(t).
\]
Using assumption~(2) together with $t \leadsto s$, we obtain
\[
P_S(s).
\]
Finally, by the equivalence between $P_S$ and $P$ on the subtype,
\[
P_S(s) \Longleftrightarrow P(\iota(s)),
\]
we conclude that
\[
P(x)=P(\iota(s))
\]
holds. This completes the induction.
\end{proof}

Theorem~\ref{thm:abstract-subtype-induction} isolates, in abstract form, the proof organization that recurs throughout the matrix decomposition arguments considered in this paper. The relation $\leadsto$ captures the passage from a current object to a locally transformed one; the slicing operation produces the smaller subproblem on which the recursive call is made; the implication
\[
P(\operatorname{slice}(s)) \Longrightarrow P_S(s)
\]
encodes the reconstruction of the current object from the solved subproblem; and the equivalence between $P_S$ and $P$ transfers the result back to the ambient universe. In this way, the theorem turns a sequence of local transformation, reduction, recursive construction, and global reconstruction into a single well-founded induction principle. In Figure~\ref{fig:proof-pattern}, these four ingredients appear as the four stages of the schematic proof flow.

\subsection{Two Core Difficulties in Formalization}

The proofs discussed in the previous subsection exhibit a recurrent mathematical pattern, but their formalization is not a direct transcription of textbook arguments. In the present setting, two difficulties appear repeatedly and determine the shape of the framework developed later. The first concerns recursion across changing matrix dimensions: the smaller problem produced by a reduction step is usually not an object of the same type as the original matrix. The second concerns the return from the smaller problem to the original one: once a decomposition has been constructed on a submatrix, it must be lifted back to the original dimension, and all relevant structural properties must be shown to survive this passage. In Lean, both issues become more visible than they are in ordinary mathematical writing, and both require dedicated infrastructure rather than ad hoc local arguments.

\subsubsection{Cross-Dimensional Recursion and Induction Organization}

In ordinary mathematical proofs, one often writes an $n\times n$ matrix in block form, isolates an $(n-1)\times(n-1)$ trailing block or a Schur complement, applies the same theorem to that smaller matrix, and then returns to the original one. On paper, this is naturally read as an induction on the size parameter. In Lean, however, the smaller object is not merely a matrix of smaller size; it inhabits a different type.

For instance, if the current matrix is represented as
\[
\mathrm{Matrix}\,(\mathrm{Fin}\,n)\,(\mathrm{Fin}\,n)\,R,
\]
then a trailing block or Schur complement typically has type
\[
\mathrm{Matrix}\,(\mathrm{Fin}\,(n-1))\,(\mathrm{Fin}\,(n-1))\,R.
\]
These two matrix types are not definitionally identical, and the recursive call cannot be expressed as if it were taking place in a single fixed ambient space. A Lean-style sketch already makes this visible:

\begin{lstlisting}
def A  : Matrix (Fin n) (Fin n) R := ...
def A1 : Matrix (Fin (n - 1)) (Fin (n - 1)) R 
    := trailingBlock A
\end{lstlisting}

Mathematically, $A_1$ is ``the same kind of object, only smaller''. Formally, it is an inhabitant of a different type. For that reason, a naive induction directly over matrices of one fixed dimension cannot capture the recursive pattern of these proofs.

The same issue arises even more clearly in rectangular settings, where one may pass from an $m\times n$ matrix to an $(m-1)\times(n-1)$ block, or more generally to a submatrix whose row and column index types change independently. The recursive step then changes not only the numerical size but the ambient type itself. What is needed is therefore not an induction over a single family member, but an induction principle formulated on a common universe that can contain matrices of all relevant sizes.

In Lean terms, the recursive theorem cannot simply have the form

\begin{lstlisting}
theorem main_n :
  forall A : Matrix (Fin n) (Fin n) R, P A := ...
\end{lstlisting}

and then recurse internally as if the smaller problem had the same type. The recursive call targets a theorem instance at a different type. To organize such proofs uniformly, one needs a larger ambient space in which objects of varying dimensions can be placed together, compared by a common measure, and related by reduction steps. This is why the later framework introduces a matrix universe and a well-founded induction principle on that universe: the main issue is not only that the size decreases, but that the decrease occurs across a family of different types.

\subsubsection{Block Reconstruction and Transfer of Properties}

Solving the smaller subproblem is only part of the argument. A recursive call produces a decomposition of a smaller matrix, but the theorem to be proved concerns the original one. The second difficulty is therefore to lift the recursively obtained factors back to the original dimension and to verify that the resulting enlarged objects satisfy the desired global statement.

On paper, this step is often written as a short block computation. One proves a decomposition for a trailing block or Schur complement, inserts the smaller factors into the lower-right corner of larger block matrices, fills the remaining entries using data from the local transformation step, and then checks the resulting identity. In Lean, however, each of these moves must be made explicit. A typical schematic situation looks like this:

\begin{lstlisting}
have hA1 : P A1 := ih A1 smaller
let F1' := liftFactorLeft  localData F1
let F2' := liftFactorRight localData F2
have hF1' : Structured F1' := ...
have hF2' : Structured F2' := ...
have hEq  : A = F1' * F2' := ...
\end{lstlisting}

Even at this level of abstraction, one can already distinguish several separate tasks: defining the lifted factors, proving that they preserve the relevant structural properties, and proving that their product gives back the original matrix.

The first issue is therefore one of construction. If a recursive call returns factors defined on a smaller index type, one must specify how they are embedded into larger block matrices of the original dimension. This is not a single generic coercion; it is a structured block-lifting operation that depends on the role played by the small factor inside the global decomposition. The second issue is one of properties. A factor that is triangular, unit triangular, orthogonal, invertible, or in some normal form on the smaller space must be shown to retain the corresponding property after lifting.

A further complication comes from reindexing. The smaller subproblem is often not represented using the same index type as the surrounding global statement, especially after row or column permutations. Thus one frequently has to move facts across equivalences of finite index types. In Lean-style pseudocode, one repeatedly encounters steps of the form

\begin{lstlisting}
let A' := reindex e1 e2 A
have hTri  : UpperTriangular U := ...
have hTri' : UpperTriangular (reindex e e U) := ...
\end{lstlisting}

Such steps are mathematically harmless, but formally they require explicit lemmas showing that the relevant properties are invariant under reindexing or transported in the intended way.

For this reason, the return from the recursive subproblem to the original theorem cannot be treated as a minor calculation. It requires a separate layer of infrastructure: block constructions, lifting operations, reindexing lemmas, and transport principles for structural properties. Without these tools, each concrete decomposition proof would have to rebuild the same bridge from the smaller problem back to the original matrix. This is the second major obstacle addressed by the framework introduced below.

\subsection{From Proof Pattern to Formal Framework}

The previous two subsections isolate the two ingredients from which the present work begins. On the one hand, several classical existence proofs for matrix decompositions exhibit a similar mathematical organization: a local step first reshapes the current matrix, a smaller subproblem is then extracted, the same kind of statement is established recursively on that subproblem, and the resulting data are finally lifted back to the original dimension. On the other hand, this pattern does not translate directly into Lean. The recursive step crosses matrix types of different dimensions, and the final reconstruction requires explicit block constructions, reindexing operations, and transport lemmas for the relevant structural properties.

The approach taken in this paper is to turn this observation into a formal design principle. Rather than treating each decomposition theorem as an isolated development, we separate the common proof pattern from the technical machinery needed to realize it. The pattern itself will be expressed in terms of decomposition goals, local transformations, reduction steps, recursive construction, and global reconstruction. The technical side will be handled by a common universe for matrices of varying dimensions, a well-founded induction principle on that universe, and a reusable collection of lifting and transport constructions. This separation makes it possible to organize different matrix decomposition proofs within a single formal framework.

The problem addressed in the rest of the paper is therefore not simply how to verify particular decomposition identities, but how to represent, in a reusable formal form, the passage from a local matrix operation to a smaller subproblem, from a smaller subproblem to a recursive conclusion, and from that recursive conclusion back to a global decomposition statement. The framework introduced below is designed precisely to make these transitions explicit and compositional.


% 找个合适的位置讲一下，对所有的分解都可以做
\section{A Unified Formal Framework}
\label{sec:framework}

This section presents the abstract architecture of the formal framework.  The guiding idea is to separate three logically different layers:

\begin{enumerate}
\item the \emph{decomposition goal}, namely what it means for a matrix to admit a desired factorization;
\item the \emph{local proof step}, namely how one transforms a matrix into a form suitable for recursive reduction;
\item the \emph{global induction mechanism}, namely how local progress is assembled into a theorem for all matrix sizes.
\end{enumerate}

The implementation mirrors this separation closely.  Decomposition statements are encoded as abstract schemas, local normalization steps are encoded as certified transformations, recursive descent is encoded by reduction methods, and global correctness is derived by a generic well-founded induction engine over a common matrix universe.

\subsection{Design Goals and General Principles}

The framework is built around modularity and reusability.  Instead of proving each decomposition theorem in a monolithic way, the project isolates the recurring proof pattern:
\[
\text{local transformation} \Longrightarrow \text{smaller subproblem} \Longrightarrow \text{lifting back}.
\]
This pattern appears in QR, PLU, and related decompositions, even though the concrete transformations and structural predicates differ.

Formally, the framework is designed so that the recursive engine does not depend on the specific algebraic meaning of the decomposition.  The engine only requires:

\begin{itemize}
\item a property \(P\) to be proved for all objects in a matrix universe;
\item a notion of base case;
\item a way to reduce a non-base object to a strictly smaller one;
\item a lifting principle turning a solution of the smaller problem into a solution of the original problem.
\end{itemize}

This makes the recursion layer instance-agnostic.  All decomposition-specific content is pushed into schemas, transformations, and reduction methods.

The code-level structure already reflects this design.  The abstract interfaces live in the \texttt{Abstractions} directory, reusable algebraic infrastructure lives in \texttt{Components}, and the universe-level induction drivers live in \texttt{Framework}.  A concrete instance such as QR then only needs to instantiate these generic interfaces.

\subsection{Abstract Representation of Decomposition Goals}

A decomposition problem is represented by an abstract schema.  Mathematically, such a schema consists of:

\begin{itemize}
\item a type \(F\) of candidate factors;
\item a predicate \(\mathsf{Prop}_F\) expressing the structural constraints on the factors;
\item a relation \(\mathsf{Eq}_F(A,f)\) expressing that the factors \(f \in F\) reconstruct the matrix \(A\).
\end{itemize}

Thus the associated existence statement has the form
\[
\exists f \in F,\qquad \mathsf{Prop}_F(f)\ \wedge\ \mathsf{Eq}_F(A,f).
\]

For example, in a QR decomposition over \(\mathbb{R}\), the factor type is a pair \((Q,R)\), the property requires \(Q\) to be orthogonal and \(R\) to be upper triangular, and the equation is
\[
A = QR.
\]

This abstraction is implemented by the following Lean definition.

\begin{lstlisting}
structure DecompositionSchema (m n : ℕ) (R : Type*) [CommRing R] where
  Factors : Type*
  property : Factors → Prop
  equation : Matrix (Fin m) (Fin n) R → Factors → Prop

def HasDecomposition {m n R} [CommRing R]
    (sch : DecompositionSchema m n R) (A : Matrix (Fin m) (Fin n) R) : Prop :=
  ∃ (factors : sch.Factors), sch.property factors ∧ sch.equation A factors
\end{lstlisting}

This definition is deliberately minimal.  It does not assume any particular factor profile or any particular algebraic identity.  As a result, the same framework can accommodate two-factor, three-factor, or more elaborate decomposition statements.

The QR instance is then obtained by instantiating this abstract interface.  Internally, the project defines the QR schema as follows.

\begin{lstlisting}
def QR_Schema_fin (n : ℕ) : DecompositionSchema n n ℝ where
  Factors := Matrix (Fin n) (Fin n) ℝ × Matrix (Fin n) (Fin n) ℝ
  property := fun (Q, R') => IsOrthogonalMatrix_fin (n := n) Q ∧ IsUpperTriangular R'
  equation := fun A (Q, R') => A = Q * R'

def HasQR_fin (A : Matrix (Fin n) (Fin n) ℝ) : Prop :=
  HasDecomposition (QR_Schema_fin n) A
\end{lstlisting}

Hence the formal QR existence proposition has the mathematical content
\[
\exists Q,R,\qquad Q^{\mathsf T}Q=I,\quad R\text{ upper triangular},\quad A=QR.
\]

\subsection{Abstract Representation of Local Transformations}

A local transformation is understood as a certified procedure that aims at a specific target condition.  It is not merely an endomorphism on matrices; rather, it is a package containing:

\begin{itemize}
\item a goal predicate \(\mathsf{Goal}\);
\item a type of transformation parameters;
\item an application map;
\item a constructive search procedure producing a transformation whenever the goal is not yet satisfied;
\item a proof that the resulting transformed matrix satisfies the goal.
\end{itemize}

In mathematical terms, a transformation package expresses the principle
\[
A \notin \mathsf{Goal}
\quad\Longrightarrow\quad
\exists t,\ \mathsf{Goal}(\mathsf{apply}(t,A)).
\]

The corresponding Lean interface is the following.

\begin{lstlisting}
structure Transformation (X : Type*) where
  T : Type*
  Goal : X → Prop
  [decGoal : DecidablePred Goal]
  apply : T → X → X
  find : (x : X) → (h : ¬ Goal x) → T
  find_spec : ∀ (x : X) (h : ¬ Goal x), Goal (apply (find x h) x)
\end{lstlisting}

The project also supports composition of transformations, allowing one to build a composite local step from smaller certified steps.  This is important in practice because many decomposition algorithms proceed in stages: first one establishes a pivot or normalization condition, then one clears a row or column, and only then does one recurse.

In the QR development, the local transformation is the internal Householder step.  Its target is the condition that the first column is already in the canonical form required for slicing.  The formal result reads:

\begin{lstlisting}
lemma QR_HouseholderStep_ready_for_slice
    (A : Matrix (Fin (k + 1)) (Fin (k + 1)) ℝ) :
    QR_ReadyForSlice_fin (k := k) (QR_HouseholderStep_fin (k := k) A) := by
  simpa [QR_HouseholderStep_fin, QR_ReadyForSlice_fin, QR_FirstColumnGoal_fin] using
    sliceable_after_apply (k := k) A
\end{lstlisting}

Mathematically, this lemma says that after one Householder step, the matrix has entered the \emph{sliceable region}: the first column has been normalized so that the recursive subproblem can be extracted from the trailing block.

\subsection{Abstract Representation of Reduction Methods}

A reduction method encodes the passage from the current matrix to a smaller subproblem together with the data needed to reconstruct the original matrix afterwards.  Abstractly, a reduction method consists of:

\begin{itemize}
\item a predicate \(\mathsf{IsSliceable}(A)\);
\item a slicing map
\[
A \mapsto \mathsf{slice}(A),
\]
defined when \(A\) is sliceable;
\item a reconstruction map
\[
(A,\text{context},\text{solution of slice}) \mapsto \mathsf{reconstruct}(A,\cdots),
\]
which inserts the recursive solution back into the original ambient problem;
\item a correctness law asserting that reconstruction from the original slice yields back the original matrix.
\end{itemize}

This is formalized by the following interface.

\begin{lstlisting}
structure ReductionMethod (m n slice_m slice_n : ℕ) (R : Type*) [CommRing R] where
  IsSliceable : Matrix (Fin m) (Fin n) R → Prop
  slice : (A : Matrix (Fin m) (Fin n) R) → (hA : IsSliceable A) →
    Matrix (Fin slice_m) (Fin slice_n) R
  reconstruct : (A : Matrix (Fin m) (Fin n) R) → (hA : IsSliceable A) →
                (slice_sol : Matrix (Fin slice_m) (Fin slice_n) R) →
                Matrix (Fin m) (Fin n) R
  reconstruct_slice_eq : ∀ (A : Matrix (Fin m) (Fin n) R) (hA : IsSliceable A),
                           reconstruct A hA (slice A hA) = A
\end{lstlisting}

The identity
\[
\mathsf{reconstruct}(A,h_A,\mathsf{slice}(A,h_A)) = A
\]
is the formal guarantee that the reduction really preserves all boundary data needed for lifting.

A typical example is the standard lower-right submatrix reduction, implemented in the project as \texttt{SubmatrixMethod}.  It removes the first row and first column and treats the trailing block as the recursive problem:

\begin{lstlisting}
noncomputable def SubmatrixMethod (n m : ℕ) (R : Type*) [CommRing R]
    (IsSliceable_def : Matrix (Fin (n + 1)) (Fin (m + 1)) R → Prop) :
    Abstractions.ReductionMethod (n + 1) (m + 1) n m R where
  IsSliceable := IsSliceable_def
  slice := fun A _hA ↦ A.submatrix Fin.succ Fin.succ
  reconstruct := fun A _hA slice_sol ↦
    let e_ι := finSuccEquivSum n
    let e_κ := finSuccEquivSum m
    let A' := reindex e_ι e_κ A
    let A₁₁ := A'.toBlocks₁₁
    let A₁₂ := A'.toBlocks₁₂
    let A₂₁ := A'.toBlocks₂₁
    let blocks := fromBlocks A₁₁ A₁₂ A₂₁ slice_sol
    blocks.reindex e_ι.symm e_κ.symm
\end{lstlisting}

Mathematically, this is exactly the passage
\[
\begin{pmatrix}
A_{11} & A_{12}\\
A_{21} & A_{22}
\end{pmatrix}
\longmapsto
A_{22},
\]
with reconstruction given by reinserting the recursively solved lower-right block into the original block matrix.

\subsection{Strategy Objects and Branch Composition}

A strategy combines a local transformation with a reduction method and augments them with measure-theoretic information needed for recursion.  Formally, a strategy contains:

\begin{itemize}
\item a transformation;
\item a reduction method;
\item a proof that the transformation target coincides with the reduction's sliceability predicate;
\item a measure \(\mu\) on the current problem;
\item a measure \(\mu_{\mathrm{slice}}\) on the smaller subproblem;
\item monotonicity and strict progress conditions.
\end{itemize}

The Lean definition is:

\begin{lstlisting}
structure ReductionStrategy (m n slice_m slice_n : ℕ) (R : Type*) [CommRing R] where
  transform : Transformation (Matrix (Fin m) (Fin n) R)
  reduction : ReductionMethod m n slice_m slice_n R
  goal_is_sliceable : transform.Goal = reduction.IsSliceable
  μ : Matrix (Fin m) (Fin n) R → ℕ
  μ_slice : Matrix (Fin slice_m) (Fin slice_n) R → ℕ
  μ_mono :
    ∀ (A : Matrix (Fin m) (Fin n) R) (t : transform.T),
      μ (transform.apply t A) ≤ μ A
  slice_progress :
    ∀ (A : Matrix (Fin m) (Fin n) R) (hA : reduction.IsSliceable A),
      μ_slice (reduction.slice A hA) < μ A
\end{lstlisting}

The two key inequalities are:
\[
\mu(\text{transform}(A)) \le \mu(A),
\qquad
\mu_{\mathrm{slice}}(\text{slice}(A)) < \mu(A).
\]
The first ensures that preparatory normalization does not increase complexity; the second is the actual source of recursive descent.

The project derives from a strategy a local reachability principle stating that every non-base matrix can be moved, possibly trivially, into a sliceable form whose recursive slice is strictly smaller.  This is encapsulated by \texttt{ReductionStrategy.mk\_reach}.

Branching behavior is handled compositionally.  For example, the reduction-level combinator
\texttt{ReductionMethod.try\_else} combines two reduction methods \(M_1\) and \(M_2\) into a single method that succeeds whenever either branch is sliceable:
\[
\mathsf{IsSliceable}(A) \equiv \mathsf{IsSliceable}_1(A) \lor \mathsf{IsSliceable}_2(A).
\]
This is useful for proof procedures that distinguish, for instance, between a generic pivot case and a degenerate zero-column case.

\subsection{Overall Workflow of the Framework}

The overall workflow can now be summarized mathematically as follows.  Given a decomposition problem, one first defines an abstract schema \(\mathcal{S}\).  One then constructs a local strategy \(\Sigma\) whose transformation stage sends a matrix into a sliceable form and whose reduction stage extracts a smaller subproblem.  Finally, one proves a lifting theorem showing that a decomposition of the slice implies a decomposition of the original matrix.

The global proof pattern is therefore
\[
\text{decomposition schema}
\Longrightarrow
\text{local strategy}
\Longrightarrow
\text{strictly smaller slice}
\Longrightarrow
\text{lifting}
\Longrightarrow
\forall A,\ \text{HasDecomposition}(A).
\]

At the code level, this interaction is split across the abstraction layer, the component layer, and the universe-level induction layer.  This is exactly what makes the framework reusable: once these interfaces are instantiated, the global theorem is obtained from generic infrastructure rather than from a custom recursive proof.

\section{Core Technical Components}
\label{sec:core-tech}

\subsection{The Matrix Universe}

Recursive matrix arguments change dimension, so they cannot be expressed directly inside one fixed matrix type.  The framework therefore introduces a \emph{matrix universe} that contains matrices of all sizes in a single ambient type.

For rectangular matrices the universe is
\[
\mathcal{U}_{\mathrm{rect}}(R)
=
\sum_{(m,n)\in \mathbb{N}\times\mathbb{N}}
\mathrm{Mat}_{m\times n}(R),
\]
and for square matrices the universe is
\[
\mathcal{U}_{\mathrm{sq}}(R)
=
\sum_{n\in\mathbb{N}}
\mathrm{Mat}_{n\times n}(R).
\]

This is implemented directly in Lean:

\begin{lstlisting}
structure FinRectFamily (m n : ℕ) (R : Type*) where
  A : Matrix (Fin m) (Fin n) R

abbrev FinRectUniverse (R : Type*) :=
  Σ (dims : ℕ × ℕ), FinRectFamily dims.1 dims.2 R

structure FinSqFamily (n : ℕ) (R : Type*) where
  A : Matrix (Fin n) (Fin n) R

abbrev FinSqUniverse (R : Type*) := Σ (n : ℕ), FinSqFamily n R
\end{lstlisting}

The framework also isolates the positive-dimensional region as a subtype.  For square matrices this is:
\[
\mathcal{U}_{\mathrm{sq}}^{+}(R)
=
\{x\in \mathcal{U}_{\mathrm{sq}}(R) : \dim(x)>0\}.
\]
This distinction is mathematically natural: zero-dimensional matrices are base cases, while positive-dimensional matrices are the objects on which recursive slicing is meaningful.

\subsection{Well-Founded Induction and the Global Recursion Engine}

The core induction engine works on a general universe \(X\) equipped with a well-founded relation.  Its most abstract form is a reduction-based induction principle: if every non-base object can be linked to a strictly smaller one and the desired property can be reconstructed from the smaller one, then the property holds everywhere.

The central engine is formalized as:

\begin{lstlisting}
theorem induction_by_reduction
    {rel : X → X → Prop} (hwf : WellFounded rel)
    (BaseSet : X → Prop)
    {r : X → X → Prop} {P : X → Prop}
    (h_trans : Transport r P)
    (IsReducible : X → Prop)
    (decompose : ∀ {x : X}, IsReducible x → X)
    (reconstruct : ∀ {x : X} (hx : IsReducible x), P (decompose hx) → P x)
    (prove_on_base : ∀ {x : X}, BaseSet x → P x)
    (reach_from_non_base : ∀ {x : X}, ¬ BaseSet x →
      ∃ y, ∃ (hy : IsReducible y), r y x ∧ rel (decompose hy) x)
    : ∀ (x : X), P x := by
  ...
\end{lstlisting}

Mathematically, this theorem expresses the following induction pattern:

\begin{theorem}[Reduction-based induction principle]
Let \(X\) be a type equipped with a well-founded relation \(\prec\).  Suppose \(P\) is a property on \(X\), \(B \subseteq X\) is a base set, and every \(x \notin B\) can be reached from some reducible \(y\) such that the reduced problem \(\mathrm{decompose}(y)\) satisfies
\[
\mathrm{decompose}(y) \prec x.
\]
If \(P\) holds on \(B\), if \(P\) can be reconstructed from \(\mathrm{decompose}(y)\) to \(y\), and if \(P\) is preserved along the relation \(r\), then \(P(x)\) holds for all \(x\in X\).
\end{theorem}

A particularly important specialization is the subtype-driven induction principle \texttt{induction\_on\_subtype'}.  It is the formal device that supports the project's main proof pattern: work recursively only on a distinguished subtype of positive-dimensional objects, but derive a theorem on the whole universe.

\begin{lstlisting}
theorem induction_on_subtype'
    (SubX : Type*) (toX : SubX → X)
    (μ : X → α) (relα : α → α → Prop) (hwf : WellFounded relα)
    (P : X → Prop)
    (P_sub : SubX → Prop)
    (P_compat : ∀ (x_sub : SubX), P_sub x_sub ↔ P (toX x_sub))
    (r_sub : SubX → SubX → Prop)
    (IsSliceable_sub : SubX → Prop)
    (slice_sub : ∀ (x_sub : SubX), IsSliceable_sub x_sub → X)
    ...
    : ∀ (x : X), P x := by
  ...
\end{lstlisting}

Its mathematical content is as follows.

\begin{theorem}[Subtype-driven well-founded induction]
Let \(X\) be a universe, let \(S\) be a distinguished subtype embedded into \(X\), and let \(\mu : X \to \alpha\) be a measure into a well-founded ordered set \((\alpha,\prec)\).  Assume that all nontrivial recursive steps are defined only on \(S\), that every non-base element of \(S\) can be locally transformed into a sliceable element of \(S\), and that the corresponding slice has strictly smaller measure.  Assume further that a proof on the slice can be lifted back to the sliceable object and then transported back along the local transformation.  Then the target property holds for every element of the whole universe \(X\).
\end{theorem}

This theorem is exactly what allows the project to formalize the ubiquitous pattern
\[
\text{local normalization}
\Longrightarrow
\text{recursive call on a smaller slice}
\Longrightarrow
\text{lifting back to the original matrix}.
\]

\subsection{Block Lifting Mechanisms}

Once a recursive subproblem has been solved, one must lift its factors back to a decomposition of the original matrix.  The framework provides reusable lifting cores for this purpose.  At the mathematical level, these lifting theorems operate on block matrices
\[
A =
\begin{pmatrix}
A_{11} & A_{12}\\
A_{21} & A_{22}
\end{pmatrix},
\]
assuming that the recursive factorization has already been established for the trailing block or for a derived slice such as a Schur complement.

For example, in Schur-type reductions one passes from a factorization of
\[
A_{22} - A_{21}A_{11}^{-1}A_{12}
\]
to a factorization of the whole matrix \(A\).  In zero-column reductions one passes from a factorization of the lower-right block \(A_{22}\) to a factorization of
\[
\begin{pmatrix}
0 & A_{12}\\
0 & A_{22}
\end{pmatrix}.
\]

The project implements generic lifting kernels for these patterns.  One example is the generic Schur-pattern lifting core:

\begin{lstlisting}
noncomputable def lift_schur_pattern
    ...
    (A : Matrix (Fin (k + 1)) (Fin (k + 1)) R)
    (h_pivot_unit : IsUnit (A 0 0))
    ...
    (h_slice_eq : subF₁ * ((Reductions.SchurMethod k R).slice A h_pivot_unit) = subF₂ * subF₃)
    : { res : LiftedThreeFactors PropF₁ PropF₂ PropF₃ // res.F₁ * A = res.F₂ * res.F₃ } :=
  ...
\end{lstlisting}

The important point is that this lifting theorem is generic in the structural predicates on the factors, provided those predicates are stable under the relevant block reindexings.  Hence a single lifting core can support multiple decomposition instances.

\subsection{Reindexing and the Transfer of Structural Properties}

Recursive block arguments constantly move between different index presentations of the same matrix.  For instance, a matrix in \(\mathrm{Fin}(k+1)\) is reindexed into a block world of the form \(\mathrm{Fin}(1)\oplus \mathrm{Fin}(k)\), where block extraction is more natural.  The framework therefore provides a general transport library for structural properties under reindexing.

The foundational equivalence used for block decompositions is the computational splitting
\[
\mathrm{Fin}(n+1) \simeq \mathrm{Fin}(1)\oplus \mathrm{Fin}(n),
\]
implemented as:

\begin{lstlisting}
def finSuccEquivSum (n : ℕ) : Fin (n + 1) ≃ Fin 1 ⊕ Fin n where
  toFun := Fin.cases (Sum.inl 0) (fun i => Sum.inr i)
  invFun := Sum.elim (fun _ => 0) Fin.succ
  left_inv := by
    intro x; cases x using Fin.cases <;> simp
  right_inv := by
    intro x
    rcases x with (y | i) <;> simp
    rw [Subsingleton.elim y 0]
\end{lstlisting}

This equivalence underlies block extraction identities such as the project lemma
\[
A[\mathrm{succ},\mathrm{succ}] = \text{bottom-right block of } \mathrm{reindex}(A),
\]
formalized by:

\begin{lstlisting}
lemma submatrix_succ_eq_toBlocks₂₂ {n m : ℕ} {R : Type*}
    (A : Matrix (Fin (n + 1)) (Fin (m + 1)) R) :
    A.submatrix Fin.succ Fin.succ =
      (reindex (finSuccEquivSum n) (finSuccEquivSum m) A).toBlocks₂₂ := by
  rfl
\end{lstlisting}

The reindexing library also proves that structural predicates are preserved under suitable equivalences.  For example, upper triangularity is stable under order-preserving reindexing:

\begin{lstlisting}
lemma isUpperTriangular_reindex
    (e : ι ≃ ι') (h_mono : StrictMono e) (A : Matrix ι ι R) :
    IsUpperTriangular A ↔ IsUpperTriangular (A.reindex e e) := by
  ...
\end{lstlisting}

Such lemmas are indispensable in lifting proofs, because the recursive factorization is often constructed in a block-reindexed coordinate system and must then be transported back to the original indexing.

\subsection{Typical Reduction Components}

The framework contains several reusable reduction modules representing common recursive moves in matrix decomposition proofs.

\paragraph{Submatrix reduction.}
This is the default reduction mechanism for problems where, after a suitable local normalization, one simply recurses on the lower-right block.  Mathematically, it corresponds to the operation
\[
A \mapsto A_{22}.
\]

\paragraph{Zero-column reduction.}
This reduction handles the degenerate case in which the first column has already vanished.  The slice is again the trailing block, but the reconstruction is specialized to the block form
\[
\begin{pmatrix}
0 & A_{12}\\
0 & A_{22}
\end{pmatrix}.
\]
The corresponding project code begins as follows:

\begin{lstlisting}
noncomputable def ZeroColumnMethod (n m : ℕ) (R : Type*) [CommRing R] :
    Abstractions.ReductionMethod (n + 1) (m + 1) n m R where
  IsSliceable := fun A ↦ ∀ i, A i 0 = 0
  slice := fun A _hA ↦ A.submatrix Fin.succ Fin.succ
  reconstruct := fun A hA slice_sol ↦
    ...
\end{lstlisting}

\paragraph{Schur-complement reduction.}
For square matrices over a field, the framework also provides a Schur-based reduction under the invertibility hypothesis on the leading entry:
\[
A \mapsto A_{22} - A_{21}A_{11}^{-1}A_{12}.
\]
This is implemented by:

\begin{lstlisting}
noncomputable def SchurMethod (n : ℕ) (R : Type*) [Field R] :
    Abstractions.ReductionMethod (n + 1) (n + 1) n n R where
  IsSliceable := fun A ↦ IsUnit (A 0 0)
  slice := fun A hA ↦
    let A' := reindex (finSuccEquivSum n) (finSuccEquivSum n) A
    let A₁₁ := A'.toBlocks₁₁
    let A₁₂ := A'.toBlocks₁₂
    let A₂₁ := A'.toBlocks₂₁
    let A₂₂ := A'.toBlocks₂₂
    let inv_A₀₀ : R := (IsUnit.unit hA).inv
    A₂₂ - A₂₁ * (!![inv_A₀₀]) * A₁₂
\end{lstlisting}

These components provide a concrete library of recursive descent mechanisms.  Because they all conform to the same abstract reduction interface, they can be combined with different transformation stages and different lifting theorems.

\subsection{Interfaces from the Framework to Concrete Instances}

To instantiate the framework for a concrete decomposition theorem, one must supply both mathematical data and proof-theoretic glue.

At the mathematical level, one provides:

\begin{itemize}
\item a decomposition schema;
\item a local transformation mechanism;
\item a reduction method;
\item a measure controlling recursion.
\end{itemize}

At the proof level, one provides:

\begin{itemize}
\item base-case theorems;
\item a lifting theorem from the slice to the original matrix;
\item a transport theorem showing that the desired decomposition property is preserved along the local transformation relation;
\item a strict decrease proof for the recursive slice.
\end{itemize}

In the square-universe driver, these data are assembled into a single package:

\begin{lstlisting}
structure SubtypeInductionInstance (X SubX : Type*) (toX : SubX → X) where
  μ : X → Nat
  μ_base : Nat
  P : X → Prop
  P_sub : SubX → Prop
  P_compat : ∀ (x_sub : SubX), P_sub x_sub ↔ P (toX x_sub)
  r_sub : SubX → SubX → Prop
  IsSliceable_sub : SubX → Prop
  slice_sub : ∀ (x_sub : SubX), IsSliceable_sub x_sub → X
  transport_sub :
    ∀ {x_sub y_sub}, r_sub y_sub x_sub → P_sub y_sub → P_sub x_sub
  lift_from_slice_sub :
    ∀ (x_sub : SubX) (hx : IsSliceable_sub x_sub),
      P (slice_sub x_sub hx) → P_sub x_sub
  ...
\end{lstlisting}

The global theorem is then obtained by a single call to the generic subtype-induction driver.  In other words, once the local mathematics of the instance has been packaged into this interface, the framework automatically provides the statement
\[
\forall x \in X,\qquad P(x).
\]

This is the precise sense in which the project separates local algebra from global recursion.  The instance development proves only the decomposition-specific lemmas; the framework supplies the global recursive theorem.


\section{Main Case Study: PLU Decomposition}
\label{sec:plu}

\subsection{Mathematical Statement of the PLU Decomposition}
Our main case study is the existence of a PLU decomposition for square matrices over a field.  In the mathematical form used by the project, the theorem is stated as follows: for every square matrix \(A \in R^{n \times n}\), there exist matrices \(P,L,U \in R^{n \times n}\) such that
\[
P A = L U,
\]
where \(P\) is a permutation matrix, \(L\) is unit lower triangular, and \(U\) is upper triangular.  The project deliberately uses the equation \(P A = L U\), rather than \(A = P L U\), because this form interacts cleanly with row permutations produced by pivoting steps.

At the internal formalization layer, matrices are indexed by \(\mathrm{Fin}(n)\), and the existence predicate is packaged as:

\begin{lstlisting}
def HasPLU_fin (A : Matrix (Fin n) (Fin n) R) : Prop :=
  HasDecomposition (PLU_Schema_fin n) A
\end{lstlisting}

This makes PLU a direct instance of the general decomposition framework: the entire theorem is reduced to providing a schema, a local strategy, lifting lemmas, and transport lemmas.

\subsection{Instantiation of the Schema for PLU}
For PLU, the abstract schema is instantiated with three square factors of the same size.  The factor type is
\[
(\mathrm{Matrix}\; n \times n)\times(\mathrm{Matrix}\; n \times n)\times(\mathrm{Matrix}\; n \times n),
\]
and the corresponding structural predicates express, respectively, that the first factor is a permutation matrix, the second is unit lower triangular, and the third is upper triangular.

In Lean, this internal schema is defined literally as:

\begin{lstlisting}
def PLU_Schema_fin (n : ℕ) : DecompositionSchema n n R where
  Factors := Matrix (Fin n) (Fin n) R ×
    Matrix (Fin n) (Fin n) R × Matrix (Fin n) (Fin n) R
  property := fun (P, L, U) =>
    IsPermutation P ∧ IsUnitLowerTriangular L ∧ IsUpperTriangular U
  equation := fun A (P, L, U) => P * A = L * U
\end{lstlisting}

This definition exhibits exactly the separation intended by the framework.  The schema records only the \emph{shape} of the decomposition theorem: the type of factors, the side conditions on them, and the reconstruction equation.  No recursive argument appears here.  The recursion enters later through the strategy and reduction interfaces.

\subsection{Local Transformations and Reduction Strategies for PLU}
The PLU proof strategy is based on two local branches.  If the leading entry \(A_{00}\) is invertible, the matrix is reduced by a Schur-complement step.  If the first column is entirely zero, the matrix is reduced by a zero-column step.  If neither condition initially holds, the strategy performs a row swap to move a nonzero pivot into the \((0,0)\)-position and then enters the Schur branch.

This logic is encoded by combining two reduction methods and a pivoting transformation:

\begin{lstlisting}
noncomputable def PLU_Reduction_fin (k : ℕ) :
    ReductionMethod (k + 1) (k + 1) k k R :=
  ReductionMethod.try_else (SchurMethod k R) (ZeroColumnMethod k k R)

noncomputable def PLU_Transform_fin (k : ℕ) :
    Transformation (Matrix (Fin (k + 1)) (Fin (k + 1)) R) :=
{
  T := Fin (k + 1)
  Goal := reduc.IsSliceable
  apply := fun i A => (swap R 0 i) * A
  ...
}
\end{lstlisting}

The key point is that the transformation stage is not itself a recursive reduction; it is a preparatory step that changes the matrix within the same dimension until one of the reduction predicates becomes true.  In particular, if the matrix is not already sliceable, the proof extracts an index \(i\) with \(A_{i0}\neq 0\), applies the swap matrix \(\mathrm{swap}(0,i)\), and obtains a matrix whose leading entry is a unit.  After that, recursion proceeds on a smaller \(k\times k\) problem.

\subsection{Key Lifting and Transfer Lemmas for PLU}
The recursive proof of PLU relies on two kinds of auxiliary lemmas.

First, there are \emph{lifting lemmas}, which rebuild a full decomposition of the original \((k+1)\times(k+1)\) matrix from a decomposition of the recursive slice.  Because PLU has two reduction branches, there are two branch-specific lifting statements:
\begin{lstlisting}
private lemma lift_from_slice_schur_case ...
private lemma lift_from_slice_zero_col_case ...

lemma lift_from_slice_plu_fin {k : ℕ}
    (A : Matrix (Fin (k + 1)) (Fin (k + 1)) R)
    (hA : (PLU_Reduction_fin (R := R) k).IsSliceable A)
    (h_slice : HasPLU_fin ((PLU_Reduction_fin (R := R) k).slice A hA)) :
    HasPLU_fin A := by
  ...
\end{lstlisting}
Internally, these lemmas use reusable block-lifting wrappers from the component layer, notably `lift_permutation_unitLower_upper_schur` and `lift_permutation_unitLower_upper_zero_col`, which package the generic block reconstruction arguments for the PLU factor profile.

Second, there is a \emph{transport lemma} for the transformation stage.  A pivoting step changes the matrix by left multiplication with a swap matrix, so the decomposition property must be transported backward along the strategy relation:
\begin{lstlisting}
lemma transport_plu_fin {k : ℕ}
    {A B : Matrix (Fin (k + 1)) (Fin (k + 1)) R}
    (hr : (PLU_Strategy_fin (R := R) k).r B A) (hB : HasPLU_fin B) :
    HasPLU_fin A := by
  ...
\end{lstlisting}
The proof is algebraically simple but conceptually important: if \(P B = L U\) and \(B = (\mathrm{swap}(0,t))A\), then
\[
(P \cdot \mathrm{swap}(0,t))A = L U,
\]
and the new left factor is again a permutation matrix because permutation matrices are closed under multiplication.

\subsection{The Main Formalization Theorem for PLU}
Once the schema, strategy, transport, lifting, and base case are supplied, the generic square-subtype induction framework produces the global theorem automatically.  The internal theorem exported by the PLU instance is:

\begin{lstlisting}
theorem exists_plu_decomposition_fin {n : ℕ} {R : Type*}
    [Field R] [DecidableEq R]
    (A : Matrix (Fin n) (Fin n) R) : HasPLU_fin A := by
  simpa using
    (SquareSubtypeInductionInstance.prove_for_fin
      (inst := PLU_Instance (R := R)) n A)
\end{lstlisting}

This statement is the clearest evidence that the framework is doing real work.  The theorem itself contains no decomposition-specific proof script: all PLU-specific mathematics has already been assembled into `PLU_Instance`, and the generic driver discharges the induction over the matrix dimension.  The base case is the trivial \(0\times 0\) matrix, while the inductive step is handled by the reach, transport, and lifting interfaces supplied by the PLU development.

\subsection{From \texttt{Fin\ n} to More General Finite Index Types}
The internal theorem above is stated on matrices indexed by \(\mathrm{Fin}(n)\), because that is the canonical setting for recursion, block splitting, and dimension-changing arguments.  The user-facing theorem, however, is stated for arbitrary finite index types equipped with the project's `FinEnum` structure.  The bridge is provided by reindexing along the canonical order isomorphism \(e : \iota \simeq_o \mathrm{Fin}(\mathrm{card}(\iota))\).

The main bridge lemma is:

\begin{lstlisting}
lemma hasPLU_reindex_iff
    (e : ι ≃o Fin (FinEnum.card ι)) (A : Matrix ι ι R) :
    HasPLU A ↔ HasPLU_fin (A.reindex e.toEquiv e.toEquiv) := by
  ...
\end{lstlisting}

This result depends on a library of property-preservation lemmas under reindexing, such as `isPermutation_reindex`, `isUnitLowerTriangular_reindex`, and `isUpperTriangular_reindex`.  With this bridge in place, the final external theorem is immediate:

\begin{lstlisting}
theorem exists_plu_decomposition [DecidableEq R] (A : Matrix ι ι R) :
    HasPLU A := by
  let e := orderIsoOfFinEnum ι
  rw [hasPLU_reindex_iff e]
  exact exists_plu_decomposition_fin (A.reindex e.toEquiv e.toEquiv)
\end{lstlisting}

This two-layer presentation is central to the design of the project.  The \(\mathrm{Fin}(n)\) layer is the internal proof engine; the `FinEnum` layer is the external semantic interface.

\subsection{Summary: The Role of PLU as the Main Case Study}
PLU is the main case study because it exercises almost every moving part of the framework.  It uses a nontrivial three-factor schema, a transformation stage based on row permutations, a branching reduction strategy, generic block-lifting infrastructure, and a final transfer theorem from canonical indices to arbitrary finite index types.  At the same time, the final theorem is obtained by a short appeal to the generic induction driver.

For that reason, the PLU development demonstrates more than the formalizability of a classical matrix theorem.  It shows that the framework can isolate decomposition-specific algebraic ideas from the global recursive proof pattern, while still remaining concrete enough to support realistic proofs with pivoting, block arguments, and index transport.

\section{Second Case Study: Framework Transfer}% and Comparative Analysis}
\label{sec:second-case}

\subsection{Option A: QR Decomposition}
The QR development serves as a second case study because it reuses the same framework architecture while changing the mathematical content of the decomposition rather substantially.  Instead of a permutation--unit-lower--upper factorization over an arbitrary field, the QR instance works over \(\mathbb{R}\), uses an orthogonal factor, and normalizes matrices by Householder reflections rather than by pivot permutations.

\subsubsection{Mathematical Goal and Structural Features of QR}
The target theorem is the usual real QR decomposition: for every square real matrix \(A\in\mathbb{R}^{n\times n}\), there exist matrices \(Q,R\in\mathbb{R}^{n\times n}\) such that
\[
A = Q R,
\]
where \(Q\) is orthogonal and \(R\) is upper triangular.  In the internal development, orthogonality is expressed by the equation \(Q^{\mathsf T}Q=1\).

Compared with PLU, the structural profile is different in three ways.  First, there are only two factors rather than three.  Second, the left factor satisfies an analytic condition, orthogonality, rather than a purely combinatorial condition such as being a permutation matrix.  Third, the normalization step is not based on pivot selection but on a canonical reflector that clears the first column below the diagonal.

\subsubsection{Instantiation of the Schema for QR}
The abstract schema is instantiated by taking the factors to be a pair \((Q,R)\) of square matrices.  The property component requires \(Q\) to be orthogonal and \(R\) to be upper triangular, while the equation component is the direct reconstruction identity \(A = Q R\).

The Lean definition is:

\begin{lstlisting}
def QR_Schema_fin (n : ℕ) : DecompositionSchema n n ℝ where
  Factors := Matrix (Fin n) (Fin n) ℝ × Matrix (Fin n) (Fin n) ℝ
  property := fun (Q, R') =>
    IsOrthogonalMatrix_fin (n := n) Q ∧ IsUpperTriangular R'
  equation := fun A (Q, R') => A = Q * R'
\end{lstlisting}

This highlights a useful aspect of the framework: the same abstract schema interface can support both a three-factor decomposition with left-action reconstruction (\(P A = L U\)) and a two-factor decomposition with direct reconstruction (\(A = Q R\)).

\subsubsection{Local Transformations and Reduction Methods for QR}
The QR strategy uses a single canonical local transformation: one Householder step.  For a matrix \(A\), the reflector \(H\) is chosen so that \(H A\) has a first column that vanishes below the top entry.  Once this has been achieved, the recursive subproblem is simply the lower-right \(k\times k\) submatrix.

In the code, this is organized as:

\begin{lstlisting}
noncomputable def QR_Transform_fin (k : ℕ) :
    Transformation (Matrix (Fin (k + 1)) (Fin (k + 1)) ℝ) where
  T := PUnit
  Goal := QR_ReadyForSlice_fin (k := k)
  apply := fun _ A => QR_HouseholderStep_fin (k := k) A
  ...

noncomputable def QR_Reduction_fin (k : ℕ) :
    ReductionMethod (k + 1) (k + 1) k k ℝ :=
  SubmatrixMethod k k ℝ (QR_ReadyForSlice_fin (k := k))
\end{lstlisting}

This differs from PLU in two important ways.  There is no branching between pivot and zero-column cases, because the Householder step deterministically produces a sliceable matrix.  And there is no need to search for a pivot row: the local normalization is encoded by an explicit reflector with its own algebraic correctness lemmas.

\subsubsection{Lifting and Property Transfer for QR}
The lifting side of QR is simpler in shape than that of PLU, but it requires different algebra.  Once the lower-left block has been cleared, a recursive decomposition of the trailing block is lifted by inserting the identity in the leading block of \(Q\) and reusing the already-normalized first row and first column data in \(R\).  The main lifting theorem is:

\begin{lstlisting}
lemma lift_from_slice_qr_fin
    (A : Matrix (Fin (k + 1)) (Fin (k + 1)) ℝ)
    (hA : QR_ReadyForSlice_fin (k := k) A)
    (h_slice : HasQR_fin (QR_Slice_fin (k := k) A hA)) :
    HasQR_fin A := by
  ...
\end{lstlisting}

The transport side is more characteristic of QR.  If the transformed matrix \(H A\) has a QR decomposition \(H A = Q R\), then the original matrix satisfies
\[
A = (H Q) R,
\]
because the Householder reflector is involutive and orthogonal.  This is formalized by:

\begin{lstlisting}
lemma transport_qr_fin
    (A : Matrix (Fin (k + 1)) (Fin (k + 1)) ℝ)
    (h_step : HasQR_fin (QR_HouseholderStep_fin (k := k) A)) :
    HasQR_fin A := by
  ...
\end{lstlisting}

Thus, in contrast with PLU, the transported factor is not a permutation-corrected left factor but a reflector-augmented orthogonal factor.

\subsubsection{The Main Formalization Theorem for QR}
As in the PLU case, once the QR-specific hooks are packaged into a square-subtype induction instance, the global theorem follows from the generic driver:

\begin{lstlisting}
theorem exists_qr_decomposition_fin {n : ℕ}
    (A : Matrix (Fin n) (Fin n) ℝ) : HasQR_fin A := by
  simpa using
    (SquareSubtypeInductionInstance.prove_for_fin
      (inst := QR_Instance) n A)
\end{lstlisting}

The user-facing theorem is then obtained by the same reindexing pattern:

\begin{lstlisting}
theorem exists_qr_decomposition (A : Matrix ι ι ℝ) : HasQR A := by
  let e := orderIsoOfFinEnum ι
  rw [hasQR_reindex_iff e]
  exact exists_qr_decomposition_fin (A.reindex e.toEquiv e.toEquiv)
\end{lstlisting}

Therefore the high-level proof architecture is unchanged: local transformation, recursive reduction, lifting from the slice, transport along the transformation relation, and final bridge to the external index language.

\subsubsection{Comparison with PLU: Reuse and New Ingredients}
The reused parts are substantial.  Both developments use the same decomposition-schema abstraction, the same transformation and reduction interfaces, the same square-subtype induction driver, the same dimension measure, and the same reindexing infrastructure for passing between \(\mathrm{Fin}(n+1)\) and block index types.  Both also rely on the same general principle that the instance-specific proof consists of transport, lifting, base, and reach lemmas.

The genuinely new ingredients are the QR-specific local algebra.  These include the Householder reflector construction, the proof that one step produces a sliceable matrix, the orthogonality lemmas for the reflector and the lifted \(Q\) factor, and the transport argument using \(H^{\mathsf T}H=1\).  In other words, the framework is reused almost verbatim, while the new work is concentrated exactly where it should be: in the decomposition-specific mathematics.

\subsubsection{Summary}
The QR case study shows that the framework is not specialized to elimination-style factorizations such as PLU.  It also supports decompositions driven by orthogonal normalization and analytic identities, while preserving the same proof architecture and much of the same infrastructure.  This is strong evidence that the framework captures a reusable formal pattern rather than a single hard-coded decomposition proof.

\subsection{Option B: Rank Normal Form}
\begin{itemize}
  \item This subsection will be retained if the second case study is chosen to be a rank normal form or a closely related echelon-type decomposition.
\end{itemize}

\subsubsection{Mathematical Goal and Structural Features}% of Rank Normal Form}
\begin{itemize}
  \item Introduce the mathematical target of the chosen rank normal form.
  \item Explain how its structure differs from PLU.
\end{itemize}

\subsubsection{Instantiation of the Schema for Rank Normal Form}
\begin{itemize}
  \item Explain how the decomposition schema is instantiated in this setting.
  \item Mention the relevant normal-form properties.
\end{itemize}

\subsubsection{Local Reductions and Strategy Organization}
\begin{itemize}
  \item Describe the local reductions and strategy organization used in this case study.
  \item Compare them informally with those of PLU.
\end{itemize}

\subsubsection{Lifting and Property Transfer for Rank Normal Form}
\begin{itemize}
  \item Explain the main lifting and transport lemmas needed to reconstruct the global normal form from smaller subproblems.
\end{itemize}

\subsubsection{The Main Formalization Theorem for Rank Normal Form}
\begin{itemize}
  \item State the main formal theorem for the chosen rank normal form.
  \item Indicate how it fits into the unified framework.
\end{itemize}

\subsubsection{Comparison with PLU: Reuse and Differences}
\begin{itemize}
  \item Clarify which parts of the infrastructure are shared with PLU.
  \item Clarify where the two instances diverge substantially.
\end{itemize}

\subsubsection{Summary}
\begin{itemize}
  \item Summarize what this second case study shows about the flexibility of the framework.
\end{itemize}

% 最后合成一节，Discussion (and Future Work)
\section{Discussion and Future Work}
% \section{Extensibility, Limitations, and Methodological Discussion}
\label{sec:discussion}

\subsection{Reusability of the Framework Across the Two Case Studies}
\begin{itemize}
  \item Summarize the extent to which the same abstract framework and technical infrastructure support both case studies.
\end{itemize}

\subsection{Where Different Instances Actually Differ}
\begin{itemize}
  \item Analyze more precisely which aspects of a decomposition instance are instance-specific.
  \item Clarify which aspects are genuinely generic.
\end{itemize}

\subsection{Methodological Significance}
\begin{itemize}
  \item Discuss the broader methodological implications of organizing formalization around reusable proof patterns rather than isolated theorems.
\end{itemize}

\subsection{Current Limitations and Future Directions}
\begin{itemize}
  \item Explain the present limitations of the framework.
  \item Outline possible directions for further generalization and automation.
\end{itemize}

% \section{Conclusion}
% \label{sec:conclusion}

% \subsection{Summary of the Main Results}
% \begin{itemize}
%   \item Briefly summarize the main results of the paper and the role of the two case studies.
% \end{itemize}

% \subsection{Outlook}
% \begin{itemize}
%   \item Conclude with a short discussion of future work and the broader prospects for formalized matrix theory.
% \end{itemize}

\backmatter

\bmhead{Supplementary information}

If your article has accompanying supplementary file/s please state so here. 

Authors reporting data from electrophoretic gels and blots should supply the full unprocessed scans for key as part of their Supplementary information. This may be requested by the editorial team/s if it is missing.

Please refer to Journal-level guidance for any specific requirements.

\bmhead{Acknowledgements}

Acknowledgements are not compulsory. Where included they should be brief. Grant or contribution numbers may be acknowledged.

Please refer to Journal-level guidance for any specific requirements.

\section*{Declarations}

Some journals require declarations to be submitted in a standardised format. Please check the Instructions for Authors of the journal to which you are submitting to see if you need to complete this section. If yes, your manuscript must contain the following sections under the heading `Declarations':

\begin{itemize}
\item Funding
\item Conflict of interest/Competing interests (check journal-specific guidelines for which heading to use)
\item Ethics approval and consent to participate
\item Consent for publication
\item Data availability 
\item Materials availability
\item Code availability 
\item Author contribution
\end{itemize}

\noindent
If any of the sections are not relevant to your manuscript, please include the heading and write `Not applicable' for that section. 

%%===================================================%%
%% For presentation purpose, we have included        %%
%% \bigskip command. Please ignore this.             %%
%%===================================================%%
\bigskip
\begin{flushleft}%
Editorial Policies for:

\bigskip\noindent
Springer journals and proceedings: \url{https://www.springer.com/gp/editorial-policies}

\bigskip\noindent
Nature Portfolio journals: \url{https://www.nature.com/nature-research/editorial-policies}

\bigskip\noindent
\textit{Scientific Reports}: \url{https://www.nature.com/srep/journal-policies/editorial-policies}

\bigskip\noindent
BMC journals: \url{https://www.biomedcentral.com/getpublished/editorial-policies}
\end{flushleft}

\begin{appendices}

\section{Section title of first appendix}\label{secA1}

An appendix contains supplementary information that is not an essential part of the text itself but which may be helpful in providing a more comprehensive understanding of the research problem or it is information that is too cumbersome to be included in the body of the paper.

%%=============================================%%
%% For submissions to Nature Portfolio Journals %%
%% please use the heading ``Extended Data''.   %%
%%=============================================%%

%%=============================================================%%
%% Sample for another appendix section			       %%
%%=============================================================%%

%% \section{Example of another appendix section}\label{secA2}%
%% Appendices may be used for helpful, supporting or essential material that would otherwise 
%% clutter, break up or be distracting to the text. Appendices can consist of sections, figures, 
%% tables and equations etc.

\end{appendices}

%%===========================================================================================%%
%% If you are submitting to one of the Nature Portfolio journals, using the eJP submission   %%
%% system, please include the references within the manuscript file itself. You may do this  %%
%% by copying the reference list from your .bbl file, paste it into the main manuscript .tex %%
%% file, and delete the associated \verb+\bibliography+ commands.                            %%
%%===========================================================================================%%

\bibliography{sn-bibliography}% common bib file
%% if required, the content of .bbl file can be included here once bbl is generated
%%\input sn-article.bbl

\end{document}
